
\subsection{算子 $\mathfrak r$ 的代数表述 III：配边论证}
\label{subsec:7-4-9}

下面考虑第~\ref{subsec:5-3-6}、\ref{subsec:5-3-2}
小节的情形：用彼此不同的哈密顿同痕
移动两个拉格朗日子流形。
沿用第~\ref{subsec:5-3-2} 小节的记号。
本小节的方法与第~\ref{subsec:7-4-5} 小节类似，
并采用配边论证。
令
\begin{equation}
\begin{aligned}
\mathfrak f^{(i)}_{\lambda,\lambda'}\colon
C(L^{(i)};\lambda)_
{\mathcal R^{\mathrm{sup}}_{l'(\lambda,\lambda')}}
\longrightarrow
C(L^{(i)\prime};\lambda,\lambda')_
{\mathcal R^{\mathrm{sup}}_{l'(\lambda,\lambda')}}
\end{aligned}
\tag{7.4.160}\label{eq:7-4-160}
\end{equation}
为仿照引理~\ref{lem:7-4-117}，
由命题~\ref{prop:4-6-55} 诱导的同态。
这里
$C(L^{(i)\prime};\lambda,\lambda')_
{\mathcal R^{\mathrm{sup}}_{l'(\lambda,\lambda')}}$
的定义方式，
与第~\ref{sec:7-2} 节定义
$C(L;\lambda,\lambda')$ 的方式相同。

定义式~\eqref{eq:7-4-160} 时，
使用辛微分同胚
\[
\psi^{(i)}_1\colon
(M,L^{(i)})\longrightarrow(M',L^{(i)\prime});
\]
它属于第~\ref{subsec:5-3-2} 小节所讨论的
哈密顿同痕族。
还要指出，式~\eqref{eq:7-4-160}
采用 $\operatorname{top}$ 模空间定义。

如第~\ref{subsec:7-4-5} 小节所讨论的，
式~\eqref{eq:7-4-160} 中的同态可扩张为
含幺、有隙且滤过的 $A_\infty$ 同态
\begin{equation}
\begin{aligned}
\mathfrak f^{(i)}_{\mathcal R^{\mathrm{sup}}}\colon
C(L^{(i)};\Lambda_{0,\mathrm{nov}})_
{\mathcal R^{\mathrm{sup}}}
\longrightarrow
C(L^{(i)\prime};\Lambda_{0,\mathrm{nov}})_
{\mathcal R^{\mathrm{sup}}},
\end{aligned}
\tag{7.4.161}\label{eq:7-4-161}
\end{equation}
其系数环为 $\mathcal R^{\mathrm{sup}}$。

\begin{theorem}{7.4.162}\label{thm:7-4-162}
存在位于式~\eqref{eq:7-4-161} 之上的
含幺、有隙且弱滤过的 $A_\infty$ 双模同态
\begin{equation}
\begin{aligned}
\varphi_{\mathcal R^{\mathrm{sup}}}\colon
H(L^{(1)},L^{(0)};\Lambda_{\mathrm{nov}})_
{\mathcal R^{\mathrm{sup}}}
\longrightarrow
H(L^{(1)\prime},L^{(0)\prime};\Lambda_{\mathrm{nov}})_
{\mathcal R^{\mathrm{sup}}},
\end{aligned}
\tag{7.4.163}\label{eq:7-4-163}
\end{equation}
它扩张由定理~\ref{thm:5-3-27} 得到的
弱滤过 $A_\infty$ 双模同态
\[
\varphi\colon
H(L^{(1)},L^{(0)};\Lambda_{\mathrm{nov}})
\longrightarrow
H(L^{(1)\prime},L^{(0)\prime};\Lambda_{\mathrm{nov}}),
\]
并具有下列性质：

\par\noindent\textup{(7.4.164.1)}\label{cond:7-4-164-1}\quad
$\varphi_{\mathcal R^{\mathrm{sup}}}$ 是同伦等价。

\par\noindent\textup{(7.4.164.2)}\label{cond:7-4-164-2}\quad
在同伦意义下，
$\varphi_{\mathcal R^{\mathrm{sup}}}$
不依赖于为这对哈密顿同痕
\[
\bigl(
\{\psi^{(1)}_\rho\}_\rho,
\{\psi^{(0)}_\rho\}_\rho
\bigr)
\]
选取的同伦。

\par\noindent\textup{(7.4.164.3)}\label{cond:7-4-164-3}\quad
在同伦意义下，
$\varphi_{\mathcal R^{\mathrm{sup}}}$
对哈密顿同痕的复合具有函子性。
\end{theorem}

本小节余下部分用于证明定理~\ref{thm:7-4-162}。
为此，需要把 $A_{n;K_1,K_2}$ 双模同态的概念
推广到弱滤过情形。

取满足条件~\ref{cond:3-1-6} 的幺半群
$G\subset\mathbb R\times2\mathbb Z$。
容易定义 $G$-有隙滤过 $A_{n;K_1,K_2}$ 双模，
以及它们之间的
$G$-有隙滤过 $A_{n;K_1,K_2}$ 双模同态。
研究弱滤过情形时则须稍加谨慎。
在此之前，先作如下说明。

\begin{lemma}{7.4.165}\label{lem:7-4-165}
设 $G'\supseteq G$，且二者均满足条件~\ref{cond:3-1-6}。
把第~\ref{sec:7-2} 节的构造中的 $G$ 换成 $G'$。
则所得 $A_\infty$ 结构在
$G'$-有隙同伦等价意义下不变。
类似地，构造所得的 $A_\infty$ 同态
（相应地，$A_\infty$ 同伦）
在模 $G'$-有隙同伦
（相应地，模 $G'$-有隙的同伦之同伦）意义下不变。
\end{lemma}

注意，对 $\beta\in G$，
范数 $\|\beta\|_{G'}$ 与 $\|\beta\|_G$
之间没有明显关系。
事实上，可能存在 $\beta,\beta'\in G$，使得
\[
\|\beta\|_{G'}<\|\beta'\|_{G'},
\qquad
\|\beta\|_G>\|\beta'\|_G.
\]
这正是引理~\ref{lem:7-4-165} 并非平凡结论的原因。

\begin{proof}
先取可数集 $\mathcal X^{(1)}(L)$，
并在由它生成的
$C^{(1)}(L;\Lambda_{0,\mathrm{nov}})$ 上
定义 $G$-$A_{n,K}$ 代数结构。
接着把 $\mathcal X^{(1)}(L)$
视为第零代链。
然后对 $G'$ 实施第~\ref{sec:7-2} 节的构造，
得到
\[
\mathcal X^{(2)}(L)\supset\mathcal X^{(1)}(L),
\]
使由前者生成的链复形
$C^{(2)}(L;\Lambda_{0,\mathrm{nov}})$
带有 $G'$-$A_{n',K'}$ 代数结构。
此外，可以假设包含映射
\[
C^{(1)}(L;\Lambda_{0,\mathrm{nov}})
\subset
C^{(2)}(L;\Lambda_{0,\mathrm{nov}})
\]
是 $G'$-$A_{n'',K''}$ 同态。

由构造，上述包含映射是
$G'$-$A_{n'',K''}$ 同伦等价。
又因对每个固定的 $G'$，
我们已经证明所得 $G'$-$A_{n',K'}$ 结构
不依赖于各项选择，引理得证。
\end{proof}

由此，在构造过程中可以自由改变 $G$。

现在回到弱滤过 $A_\infty$ 双模同态。
取满足条件~\ref{cond:3-1-6} 的 $G$。
设 $C_1,C_0$ 分别为
$G$-$A_{n,K_1}$、$G$-$A_{n,K_0}$ 代数。
仿照定义~\ref{def:7-2-63}，
在
$G\times\mathbb Z_{\geq0}\times\mathbb Z_{\geq0}$
上定义序关系 $\lesssim$。
设 $D$ 为带上边界算子
$\mathfrak n_{0,0,\beta_0}$ 的
滤过 $\Lambda_{0,\mathrm{nov}}$-模。

\begin{definition}{7.4.166}\label{def:7-4-166}
$D$ 上位于 $C_1,C_0$ 之上的
$G$-$A_{n;K_1,K_0}$ 双模结构\liftIndex{I337}，
是运算族
\[
\mathfrak n_{k_1,k_0;\beta}\colon
B_{k_1}(\overline C_1[1])
\otimes\overline D[1]
\otimes B_{k_0}(\overline C_0[1])
\longrightarrow\overline D[1],
\]
其中 $(\beta;k_1,k_0)\lesssim(n;K_1,K_0)$，
并且对每个这样的 $(\beta;k_1,k_0)$，都有
\begin{equation}
\begin{aligned}
&\sum_{\substack{
k_1^{(1)}+k_1^{(2)}=k_1+1,\;
k_0^{(1)}+k_0^{(2)}=k_0+1\\
\beta^{(1)}+\beta^{(2)}=\beta}}
\widehat{\mathfrak n}_{
k_1^{(1)},k_0^{(1)};\beta^{(1)}}
\circ
\widehat{\mathfrak n}_{
k_1^{(2)},k_0^{(2)};\beta^{(2)}}
=0.
\end{aligned}
\tag{7.4.167}\label{eq:7-4-167}
\end{equation}
\end{definition}

可以按显然方式定义两个
$G$-$A_{n;K_1,K_0}$ 双模之间的
$G$-$A_{n;K_1,K_0}$ 双模同态。
不过，本小节需要的是弱滤过
$A_\infty$ 双模同态（定义~\ref{def:5-2-1}）。

弱滤过情形稍微复杂一些。
设 $G'$ 是一个 $G$-集合
（见定义~\ref{def:5-2-3}）。
考虑分别位于
$(C_1,C_0)$ 与 $(C'_1,C'_0)$ 之上的两个
$G$-$A_{n;K_1,K_0}$ 双模 $D,D'$；
下面定义二者之间的弱滤过
$G'$-$A_{n;K_1,K_0}$ 双模同态。

\begin{definition}{7.4.168}\label{def:7-4-168}
对 $\beta\in G'$，置
\[
\|\beta\|
=
\sup\left\{
n\,\middle|\,
\beta=\beta'+\sum_{i=1}^{n}\beta_i,
\quad
\beta'\in G',\ \beta_i\in G
\right\}
+[\omega(\beta)+c].
\]
这里 $c$ 是条件 (5.2.4.1) 中定义的能量损失，
而 $[x]$ 是满足 $[x]\leq x$ 的最大整数。
\end{definition}

\begin{definition}{7.4.169}\label{def:7-4-169}
设
\[
\mathfrak f^{(i)}\colon C_i\longrightarrow C'_i
\qquad(i=1,0)
\]
是滤过 $G$-$A_{n,K_i}$ 代数同态。
位于 $\mathfrak f^{(i)}$ 之上、从 $D$ 到 $D'$ 的
一个弱滤过 $G'$-$A_{n;K_1,K_0}$ 双模同态\liftIndex{I379}，
是同态序列
\[
\varphi_{k_1,k_0,\beta}\colon
B_{k_1}(\overline C_1[1])
\otimes\overline D[1]
\otimes B_{k_0}(\overline C_0[1])
\longrightarrow\overline D'[1],
\]
其中
$(\beta;k_1,k_0)\lesssim(n;K_1,K_0)$、
$\beta\in G'$，并且对每个这样的指标都有
\begin{equation}
\begin{aligned}
&\sum_{k_1^{(1)}+k_1^{(2)}=k_1+1}
\sum_{k_0^{(1)}+k_0^{(2)}=k_0+1}
\sum_{\beta^{(1)}+\beta^{(2)}=\beta}
\Bigl(
\widehat{\mathfrak n}_{
k_1^{(1)},k_0^{(1)};\beta^{(1)}}
\circ
\widehat\varphi_{
k_1^{(2)},k_0^{(2)};\beta^{(2)}}
\notag\\[-.2em]
&\hspace{10em}
-
\widehat\varphi_{
k_1^{(2)},k_0^{(2)};\beta^{(2)}}
\circ
\widehat{\mathfrak n}_{
k_1^{(1)},k_0^{(1)};\beta^{(1)}}
\Bigr)
=0.
\end{aligned}
\tag{7.4.170}\label{eq:7-4-170}
\end{equation}
\translatornote{原书陈述式的 $\widehat{\mathfrak n}$
均无撇号；若按型别，第一项作用于 $D'$ 一侧，
似应取带撇的运算。源书如此，译文照录。见疑误 A-0741。}
式~\eqref{eq:7-4-170} 中，
$\beta^{(1)}\in G$、$\beta^{(2)}\in G'$。
\end{definition}

现在陈述定理~\ref{thm:7-4-162} 的
$A_{n;K_1,K_2}$ 版本。
考虑与该定理相同的情形。
回忆：在定理~\ref{thm:F} 的证明过程中，
我们构造了位于
$C(L^{(1)};\lambda)$--$C(L^{(0)};\lambda)$
之上的 $G$-$A_{n(\lambda);K_1(\lambda),K_2(\lambda)}$ 双模
\[
C(L^{(1)},L^{(0)};\lambda).
\]
此外，第~\ref{subsec:7-4-5} 小节还定义了
\[
C(L^{(i)};\lambda)_
{\mathcal R^{\mathrm{sup}}_{l(\lambda)}}.
\]

\begin{proposition}{7.4.171}\label{prop:7-4-171}
存在数列
\[
l'(\lambda,\lambda')\nearrow\infty,
\qquad
n'(\lambda,\lambda')\nearrow\infty,
\qquad
K'(\lambda,\lambda')\nearrow\infty,
\]
使得对任意
$l'(\lambda,\lambda')$、
$n'(\lambda,\lambda')$ 与
$K'(\lambda,\lambda')$，
都存在一个含幺滤过双模
\translatornote{源书此处写 module；但该对象同时位于紧接着给出的左右两个代数之上，后文的同态与包含映射也都明确称为 bimodule homomorphism，故译为“双模”。见疑误 A-0601。}，
其截断型别为
\[
G\text{-}A_{n'(\lambda,\lambda');
K'(\lambda,\lambda'),K'(\lambda,\lambda')}
\]
而该双模是
\begin{equation}
C(L^{(1)\prime},L^{(0)\prime};\lambda,\lambda')_
{\mathcal R^{\mathrm{sup}}_{l'(\lambda,\lambda')}}.
\tag{7.4.172}\label{eq:7-4-172}
\end{equation}
它位于
\[
C(L^{(1)\prime};\lambda,\lambda')_
{\mathcal R^{\mathrm{sup}}_{l'(\lambda,\lambda')}}
\;--\;
C(L^{(0)\prime};\lambda,\lambda')_
{\mathcal R^{\mathrm{sup}}_{l'(\lambda,\lambda')}}
\]
之上。
而且，存在位于式~\eqref{eq:7-4-160} 中
$\mathfrak f^{(i)}_{\lambda,\lambda';
\mathcal R^{\mathrm{sup}}}$ 之上的
含幺弱滤过
$G'$-$A_{n'(\lambda,\lambda');
K'(\lambda,\lambda'),K'(\lambda,\lambda')}$ 双模同态
\begin{equation}
\begin{aligned}
\varphi_{\lambda,\lambda';\mathcal R^{\mathrm{sup}}}\colon
C(L^{(1)},L^{(0)};\lambda)_
{\mathcal R^{\mathrm{sup}}_{l'(\lambda,\lambda')}}
\longrightarrow
C(L^{(1)\prime},L^{(0)\prime};\lambda,\lambda')_
{\mathcal R^{\mathrm{sup}}_{l'(\lambda,\lambda')}}.
\end{aligned}
\tag{7.4.173}\label{eq:7-4-173}
\end{equation}
此外，有包含映射
\[
\begin{aligned}
C(L^{(1)\prime},L^{(0)\prime};\lambda')_
{\mathcal R^{\mathrm{sup}}_{l'(\lambda,\lambda')}}
\subset
C(L^{(1)\prime},L^{(0)\prime};\lambda,\lambda')_
{\mathcal R^{\mathrm{sup}}_{l'(\lambda,\lambda')}}.
\end{aligned}
\]
该包含映射诱导一个
$G$-$A_{n'(\lambda,\lambda');
K'(\lambda,\lambda'),K'(\lambda,\lambda')}$ 双模同态，
并且是同伦等价。
\end{proposition}

\begin{proof}
考虑第~\ref{subsec:5-3-2} 小节式~\eqref{eq:5-3-19}
中的模空间
\begin{equation}
\begin{aligned}
\mathcal M^\circ_{k_1,k_0}
\bigl(&
(L^{(1)},\psi^{(1)}),
(L^{(0)},\psi^{(0)});
\{J_{\tau,t}\}_{\tau,t};\\[-.2em]
&[h,w],[h',w'];\operatorname{top}(\tau)
\bigr).
\end{aligned}
\tag{7.4.174}\label{eq:7-4-174}
\end{equation}
仿照第~\ref{subsec:5-3-2} 小节将它紧化，
而不采用第~\ref{subsec:7-2-14} 小节
式~\eqref{eq:7-2-308} 的方式。
再以此紧化模空间给出的对应，定义
$G'$-$A_{n'(\lambda,\lambda');
K'(\lambda,\lambda'),K'(\lambda,\lambda')}$ 双模同态
\begin{equation}
\varphi_{\lambda,\lambda'}\colon
C(L^{(1)},L^{(0)};\lambda)
\longrightarrow
C(L^{(1)\prime},L^{(0)\prime};\lambda,\lambda').
\tag{7.4.175}\label{eq:7-4-175}
\end{equation}

接着在模空间~\eqref{eq:7-4-174}
上加入 $\ell$ 个内部标记点，
并以类似方式紧化，得到
\begin{equation}
\begin{aligned}
\mathcal M_{k_1,k_0;\ell}
\bigl(&
(L^{(1)},\psi^{(1)}),
(L^{(0)},\psi^{(0)});
\{J_{\tau,t}\}_{\tau,t};\\[-.2em]
&[h,w],[h',w'];\operatorname{top}(\tau)
\bigr).
\end{aligned}
\tag{7.4.176}\label{eq:7-4-176}
\end{equation}
仿照第~\ref{subsec:7-4-5} 小节使用它，
便得到具有所需性质的式~\eqref{eq:7-4-173}。

注意，式~\eqref{eq:7-4-172} 并不是由
式~\eqref{eq:7-4-175} 中的
$C(L^{(1)\prime},L^{(0)\prime};\lambda,\lambda')$
通过下式得到的：
\[
\begin{aligned}
C(L^{(1)\prime},L^{(0)\prime};\lambda,\lambda')_
{\mathcal R^{\mathrm{sup}}_{l'(\lambda,\lambda')}}
\cong
C(L^{(1)\prime},L^{(0)\prime};\lambda,\lambda')
\otimes_{\Lambda_{0,\mathrm{nov}}}
\mathcal R^{\mathrm{sup}}_{l'(\lambda,\lambda')}.
\end{aligned}
\]
换言之，还须加入更多的链。
这个过程已反复实施，故从略。

还要指出，为定义式~\eqref{eq:7-4-173}，
先取
$\varphi^{\mathrm{geo}}_{
\lambda,\lambda';\mathcal R^{\mathrm{sup}}}$，
再用同调代数得到
$\varphi_{
\lambda,\lambda';\mathcal R^{\mathrm{sup}}}$。
采用定理~\ref{thm:7-4-149} 的证明即可做到这一点。
这一过程也已反复实施，故从略。
命题~\ref{prop:7-4-171} 得证。
\end{proof}

\begin{proof}[定理~\ref{thm:7-4-162} 的证明]
仿照第~\ref{sec:7-2} 节使用命题~\ref{prop:7-4-171}，
即可得到滤过 $A_\infty$ 双模同态
\eqref{eq:7-4-163}。
仿照第~\ref{subsec:5-3-3}、\ref{subsec:7-2-10} 小节，
可证明它在同伦意义下不依赖于各项选择。
复合哈密顿同痕时，
所得滤过 $A_\infty$ 双模同态
\eqref{eq:7-4-163} 也相应复合；
这一事实的证明与第~\ref{subsec:5-3-4} 小节的论证相同。

最后仿照第~\ref{subsec:5-3-2} 小节使用这些结论，
即可完成定理~\ref{thm:7-4-162} 的证明。
\end{proof}

\begin{remark}{7.4.177}\label{rem:7-4-177}
在当前情形下，分叉法似乎难以应用。
因此，我们不知道
\[
\operatorname{Der}B
\bigl(C_{\mathrm{para}}(L^{(1)},L^{(0)})[1]\bigr)
\otimes_{\Lambda_{0,\mathrm{nov}}}\Lambda_{\mathrm{nov}}
\]
与
\[
\operatorname{Der}B
\bigl(C_{\mathrm{para}}(L^{(1)\prime},L^{(0)\prime})[1]\bigr)
\otimes_{\Lambda_{0,\mathrm{nov}}}\Lambda_{\mathrm{nov}}
\]
在定理~\ref{thm:7-4-162} 的情形下
是否（弱）同伦等价。
第~\ref{subsec:7-2-14} 小节给出的论证
混合了两种方法，但似乎也不能证明这一结论。

尝试证明该陈述时，一个可以预见的困难如下。
在处理（有隙）滤过情形的
定理~\ref{thm:7-4-149} 的证明中，
我们实质性地使用了典范模型。
典范双模在滤过 $A_\infty$ 双模的
滤过同伦等价下得以保持，
但在弱滤过同伦等价下并不保持。
\end{remark}
